\documentclass[11pt]{article}
\usepackage[a4paper,margin=24mm]{geometry}
\usepackage{amsmath,amssymb}
\setlength{\parindent}{0pt}
\setlength{\parskip}{6pt}
\newcommand{\E}{\mathrm e}
\newcommand{\B}{B}
\begin{document}
\begin{center}
{\Large\bfseries A singleton-event counterexample to the stated\\[3pt]
quadratic divergence law}\\[7pt]
{\small Conjecture 00000007862 --- literal displayed-bound clause}
\end{center}

\textbf{Scope.} The bilingual source gives the bound
\[
 \B(\Phi)=\frac{\sum_i p_i}
 {(1-\E p^*)\prod_i(1-p_i)}
\]
and then states that \emph{the bound} diverges quadratically as
\(\E p^*\to1\). We disprove this unqualified law on a fixed admissible
family. The proof concerns the displayed bound, without assuming or
formalizing a formula for the expected restart count \(T(\Phi)\).

\textbf{Concrete probability and event model.} There is one variable
with two possible outcomes, \(0\) and \(1\). For \(0<p<1\), put
\[
 \mathbb P(0)=p,\qquad \mathbb P(1)=1-p.
\]
These masses are nonnegative and sum to one. There is one bad event,
\(A=\{0\}\). Its probability, computed by summing the masses of its
outcomes, is \(\mathbb P(A)=p\). Thus the largest event probability is
actually \(p^*=p\), attained by the sole event.

The event depends on the sole variable. Define dependency adjacency by
requiring distinct events whose variable supports overlap. This graph
has no edges and degree \(d=0\): a singleton set of events contains no
distinct pair. The usual symmetric LLL admissibility condition
\(\E p^*(d+1)<1\) therefore reduces to \(\E p<1\). Any requirement of
independence for distinct nonadjacent events is vacuous in this genuine
singleton event family.

\textbf{Approach to the critical value.} Let \(\E=\exp(1)>1\), and for
every integer \(n\geq0\) choose
\[
 p_n=\frac{1-1/(n+2)}{\E},\qquad
 \delta_n=1-\E p_n=\frac1{n+2}.
\]
Then \(0<p_n<1/\E<1\), \(\E p_n<1\), and
\(\E p_n\to1\) from below. This is an admissible sequence of actual
probability laws on the same variable-event hypergraph.

For this family the finite sum and product in the source reduce to
\(\sum_i p_i=p_n\) and \(\prod_i(1-p_i)=1-p_n\). Consequently its
\emph{actual displayed bound} is the positive number
\[
 \B_n=\frac{p_n}{\delta_n(1-p_n)}.
\]
No scalar probability or maximum is added as an unproved premise:
both are computed from the concrete event family.

\newpage
\begin{center}
{\large\bfseries The asymptotic contradiction and its formal scope}
\end{center}

\textbf{First-order behavior.} Multiplication by the critical gap gives
\[
 \delta_n\B_n=\frac{p_n}{1-p_n}
   \longrightarrow
   \frac{1/\E}{1-1/\E}
   =\frac1{\E-1}>0.
\]
Hence \(\B_n\sim(\E-1)^{-1}\delta_n^{-1}\). The displayed bound really
diverges along this family, but its pole has order one.

\textbf{Quadratic normalization.} A second multiplication gives
\[
 \delta_n^2\B_n
   =\delta_n\frac{p_n}{1-p_n}
   \longrightarrow0.
\]
For any fixed \(c>0\), convergence to zero implies that eventually
\(\delta_n^2\B_n<c\). Since \(\delta_n>0\), it follows that
\[
 \neg\bigl(\text{eventually }c\,\delta_n^{-2}\leq\B_n\bigr).
\]
Thus the bound cannot have a positive quadratic asymptotic coefficient.
It also cannot satisfy a two-sided \(\Theta(\delta_n^{-2})\) law, because
such a law would supply the positive eventual lower comparison just
excluded. This refutes the asserted universal quadratic divergence
conjunct of the source.

\textbf{What is, and is not, concluded.} The source does not specify a
worst-case supremum over instances or a growing family of event counts.
An expression can have different behavior if its remaining factors
also vary singularly in a specially chosen growing family. Such an
existential or worst-case statement is a different claim, and is not
disproved here. Our counterexample applies to the unqualified law when
it is asserted for admissible families. This interpretation is explicit
and must receive independent source-bound review.

The source's other claims concern an upper bound for \(T\), a branching
process explanation, and an adaptive-selection exponent. We neither
prove nor individually disprove those claims. A false necessary
conjunct suffices to reject the original conjunction under the stated
universal reading. No resampling algorithm is silently supplied to
repair the source's omitted definition of restart rounds.

\textbf{Lean correspondence.} \texttt{Main.lean} defines a finite
probability law with nonnegative masses and total mass one;
\texttt{instanceLaw}, \texttt{variableValue}, and \texttt{eventFamily}
give the actual variable and event model.
\texttt{pStar\_largest} proves the maximum by
\texttt{IsGreatest}; \texttt{dependencyDegree\_zero} and
\texttt{instance\_admissible} establish degree-zero admissibility.
\texttt{actualBound} retains the original finite expression.
\texttt{actual\_scaled\_square\_tendsto\_zero},
\texttt{actual\_scaled\_linear\_tendsto}, and
\texttt{actual\_no\_positive\_quadratic\_lower\_bound} prove the limits
and contradiction for this actual bound. The premise-free theorem
\texttt{literal\_bound\_counterexample} collects the certificate.

\textit{Verification status.} The pinned Lean package and this proof
are author artifacts. Deterministic replay, independent semantic
review, and acceptance belong to the manager; this document does not
self-certify them.
\end{document}
